\section{Jailbreak Attacks}
\label{sec:jailbreak-attacks}

Chương này hệ thống hóa các nhóm phương pháp jailbreak đối với mô hình
ngôn ngữ lớn đã được căn chỉnh an toàn. Mục tiêu của chương không phải là
cung cấp hướng dẫn tạo nội dung gây hại, mà là phân tích các cơ chế tấn
công ở mức khái niệm để phục vụ đánh giá an toàn và thiết kế phòng thủ.
Các phương pháp được phân loại theo quyền truy cập của attacker, số lượt
tương tác, mức độ tự động hóa, khả năng chuyển giao và khả năng thích nghi
trước defense.

\subsection{Taxonomy các phương pháp jailbreak}
\label{subsec:jailbreak-taxonomy}

Jailbreak attack là quá trình xây dựng hoặc tối ưu đầu vào nhằm khiến mô
hình tạo ra phản hồi không phù hợp với chính sách an toàn. Xét một yêu cầu
gốc $x$ thuộc nhóm không được phép, attacker tìm một biến đổi
$\mathcal{T}$ để tạo ra đầu vào mới

\begin{equation}
    x' = \mathcal{T}(x),
    \label{eq:jailbreak-transformed-input}
\end{equation}

trong đó $x'$ vẫn giữ mục tiêu ngữ nghĩa của $x$ nhưng làm suy yếu hành vi
từ chối của mô hình. Mục tiêu khái quát của attack có thể được biểu diễn là

\begin{equation}
    F_{\theta}(x', s, c)
    \notin
    \mathcal{Y}_{\mathrm{safe}},
    \label{eq:jailbreak-attack-goal}
\end{equation}

trong đó $F_{\theta}$ là mô hình, $s$ là system instruction, $c$ là ngữ
cảnh bổ sung và $\mathcal{Y}_{\mathrm{safe}}$ là tập phản hồi phù hợp với
chính sách an toàn.

Chương~\ref{sec:background} đã trình bày threat model của jailbreak theo
các chiều như quyền truy cập, số lượt tương tác, ngân sách truy vấn, mức
độ hiểu biết của attacker và khả năng thích nghi. Dựa trên nền tảng đó,
chương này tập trung vào cách các jailbreak attacks được xây dựng trong
thực tế, tức là phân loại theo cơ chế tạo hoặc tối ưu đầu vào.

Bảng~\ref{tab:jailbreak-method-families} tóm tắt các nhóm phương pháp
jailbreak được phân tích trong chương này.

\begin{table}[H]
\centering
\caption{Các nhóm phương pháp jailbreak được phân tích trong chương này.}
\label{tab:jailbreak-method-families}
\begin{tabularx}{\textwidth}{
    >{\bfseries}p{0.24\textwidth}
    p{0.30\textwidth}
    X
}
\toprule
Nhóm phương pháp
& Cơ chế chính
& Đặc điểm tiêu biểu \\
\midrule

Manual jailbreak
& Viết lại yêu cầu bằng bối cảnh, vai trò hoặc cách diễn đạt khác
& Đơn giản, không cần truy cập nội bộ mô hình, nhưng phụ thuộc nhiều vào
kinh nghiệm của người thiết kế. \\

Template-based attack
& Sử dụng một mẫu prompt cố định rồi thay thế nội dung mục tiêu
& Dễ mở rộng trên nhiều yêu cầu, nhưng dễ bị nhận diện nếu template quá
cố định. \\

Obfuscation và encoding
& Biến đổi hình thức bề mặt của đầu vào nhưng vẫn giữ mục tiêu ngữ nghĩa
& Khai thác hạn chế của bộ lọc từ khóa, classifier hoặc các mẫu từ chối
dựa trên đặc trưng bề mặt. \\

Multi-turn jailbreak
& Dẫn dắt mô hình qua nhiều lượt hội thoại thay vì yêu cầu trực tiếp
& Rủi ro tích lũy theo ngữ cảnh, khó đánh giá nếu chỉ xét từng lượt riêng
lẻ. \\

White-box optimization
& Tối ưu chuỗi token hoặc adversarial suffix dựa trên thông tin nội bộ
của mô hình
& Hữu ích để nghiên cứu giới hạn của safety alignment, nhưng giả định
quyền truy cập mạnh. \\

Black-box automated jailbreak
& Tự động sinh, đánh giá và tinh chỉnh prompt thông qua phản hồi của
target model
& Phù hợp với bối cảnh dịch vụ đóng, nhưng phụ thuộc vào ngân sách truy
vấn và chất lượng evaluator. \\

Transferable attack
& Tạo attack trên một mô hình hoặc tập yêu cầu rồi chuyển sang mục tiêu
khác
& Cho thấy các mô hình khác nhau có thể chia sẻ điểm yếu trong cơ chế
từ chối. \\

Adaptive attack
& Điều chỉnh chiến lược dựa trên defense hoặc phản hồi của hệ thống
& Quan trọng khi đánh giá độ bền của defense trong điều kiện attacker có
khả năng thích nghi. \\

\bottomrule
\end{tabularx}
\end{table}

Taxonomy này cho thấy jailbreak không phải một kỹ thuật đơn lẻ mà là một
nhóm rộng các chiến lược tấn công. Một attack có thể đồng thời thuộc nhiều
nhóm, chẳng hạn black-box, multi-turn, automated và adaptive.

\subsection{Manual và template-based attacks}
\label{subsec:manual-template-attacks}

Các jailbreak đầu tiên thường được xây dựng thủ công. Attacker viết lại
yêu cầu ban đầu bằng cách thay đổi vai trò, bối cảnh hoặc hình thức chỉ
dẫn nhằm làm mô hình hiểu sai ý định an toàn. Những attack này thường dựa
trên khả năng khai thác sự không nhất quán giữa helpfulness và refusal:
mô hình được huấn luyện để hỗ trợ người dùng, nhưng đồng thời phải tuân thủ
giới hạn an toàn.

Manual jailbreak có thể sử dụng nhiều chiến lược ở mức khái niệm, chẳng
hạn:

\begin{itemize}[leftmargin=1.2cm]
    \item thay đổi bối cảnh để làm yêu cầu trông giống một tình huống giả
    định hoặc phân tích học thuật;
    \item yêu cầu mô hình nhập vai vào một vai trò không bị ràng buộc bởi
    chính sách thông thường;
    \item chia nhỏ mục tiêu thành nhiều phần để làm giảm khả năng phát hiện
    của bộ lọc;
    \item diễn đạt lại yêu cầu sao cho mô hình tập trung vào tính hữu ích
    thay vì nhận diện rủi ro.
\end{itemize}

Template-based attack tổng quát hóa manual attack bằng cách sử dụng một
mẫu prompt có thể áp dụng cho nhiều yêu cầu khác nhau. Thay vì thiết kế
lại toàn bộ prompt cho từng mục tiêu, attacker chỉ thay thế một số trường
trong template, chẳng hạn loại nhiệm vụ, bối cảnh hoặc nội dung cần kiểm
thử. Cách này giúp mở rộng attack nhanh hơn nhưng thường kém bền vững, vì
template cố định dễ bị nhận diện bởi các mô hình hoặc bộ lọc an toàn.

Ưu điểm của manual và template-based attacks là đơn giản, không yêu cầu
truy cập vào mô hình và dễ triển khai trong black-box setting. Tuy nhiên,
chúng thường phụ thuộc nhiều vào kinh nghiệm của người thiết kế, dễ bị lỗi
thời khi mô hình được cập nhật và có khả năng chuyển giao không ổn định.

\subsection{Obfuscation và encoding attacks}
\label{subsec:obfuscation-encoding-attacks}

Obfuscation attack tìm cách thay đổi hình thức bề mặt của đầu vào mà vẫn
giữ lại mục tiêu ngữ nghĩa. Ý tưởng chính là làm cho yêu cầu trở nên khó
nhận diện hơn đối với các bộ lọc dựa trên từ khóa, classifier hoặc các mẫu
từ chối đã học trong quá trình safety alignment.

Các biến đổi thường gặp ở mức khái niệm bao gồm:

\begin{itemize}[leftmargin=1.2cm]
    \item thay đổi chính tả, dấu câu hoặc cách phân tách từ;
    \item dùng ký hiệu, mã hóa hoặc biểu diễn trung gian;
    \item diễn đạt vòng vo thay vì nêu trực tiếp mục tiêu;
    \item sử dụng nhiều ngôn ngữ hoặc trộn các hệ chữ viết;
    \item tách một yêu cầu thành nhiều mảnh ngữ nghĩa rời rạc.
\end{itemize}

Những attack này khai thác khoảng cách giữa cách con người hiểu nội dung
và cách mô hình hoặc guard model biểu diễn nội dung dưới dạng token. Nếu
bộ lọc an toàn phụ thuộc quá nhiều vào đặc trưng bề mặt, một biến đổi nhỏ
có thể làm giảm khả năng phát hiện. Tuy nhiên, obfuscation không bảo đảm
thành công nếu hệ thống có cơ chế chuẩn hóa đầu vào, semantic filtering
hoặc đánh giá dựa trên toàn bộ ngữ cảnh.

Trong báo cáo này, obfuscation được xem là một nhóm attack quan trọng vì
nó cho thấy phòng thủ ở mức từ khóa hoặc mẫu cố định thường không đủ. Một
defense hiệu quả cần đánh giá ý định ngữ nghĩa của yêu cầu, không chỉ dựa
trên hình thức bề mặt của prompt.

\subsection{Multi-turn jailbreak}
\label{subsec:multi-turn-jailbreak}

Single-turn jailbreak cố gắng vượt qua cơ chế an toàn trong một truy vấn
duy nhất. Ngược lại, multi-turn jailbreak xây dựng attack qua nhiều lượt
hội thoại. Thay vì yêu cầu trực tiếp mục tiêu cuối cùng, attacker có thể
dẫn dắt mô hình qua một chuỗi câu hỏi trung gian, làm cho ngữ cảnh hội
thoại dần dịch chuyển về phía hành vi không mong muốn.

Một hội thoại multi-turn có thể được biểu diễn bởi chuỗi

\begin{equation}
    h_t = (u_1, y_1, u_2, y_2, \ldots, u_t),
    \label{eq:multi-turn-history}
\end{equation}

trong đó $u_i$ là lượt người dùng và $y_i$ là phản hồi của mô hình. Ở mỗi
lượt, phản hồi tiếp theo phụ thuộc không chỉ vào yêu cầu hiện tại mà còn
vào lịch sử hội thoại:

\begin{equation}
    y_t = F_{\theta}(h_t, s, c).
    \label{eq:multi-turn-generation}
\end{equation}

Multi-turn jailbreak khó đánh giá hơn single-turn vì rủi ro có thể không
nằm trong từng lượt riêng lẻ mà xuất hiện từ sự tích lũy ngữ cảnh. Một câu
hỏi trung gian có thể trông hợp lệ khi xét độc lập, nhưng khi kết hợp với
các lượt trước đó lại phục vụ cho mục tiêu không an toàn.

Nhóm attack này đặc biệt quan trọng đối với agentic systems. Agent thường
có bộ nhớ, kế hoạch và trạng thái trung gian; do đó attacker có thể tác
động từng bước đến quá trình ra quyết định thay vì chỉ tấn công phản hồi
cuối cùng. Đây là lý do evaluation cho agent cần xem xét toàn bộ trajectory
thay vì chỉ kiểm tra output cuối.

\subsection{White-box optimization}
\label{subsec:white-box-optimization}

Trong thiết lập white-box, attacker có quyền truy cập vào một phần thông
tin nội bộ của mô hình, chẳng hạn tham số, gradient hoặc logit. Khi đó,
jailbreak có thể được mô hình hóa như một bài toán tối ưu đầu vào. Một
hướng tiêu biểu là tìm adversarial suffix $a$ để nối vào yêu cầu gốc $x$:

\begin{equation}
    x' = x \oplus a,
    \label{eq:adversarial-suffix-input}
\end{equation}

trong đó $\oplus$ là phép nối chuỗi. Mục tiêu là tìm suffix làm tăng xác
suất mô hình tạo phản hồi thuộc dạng mà attacker mong muốn. Bài toán tổng
quát có thể viết dưới dạng

\begin{equation}
    a^{*}
    =
    \arg\min_{a \in \mathcal{V}^{m}}
    \mathcal{L}
    \left(
        F_{\theta}(x \oplus a)
    \right),
    \label{eq:adversarial-suffix-optimization}
\end{equation}

trong đó $\mathcal{V}$ là vocabulary, $m$ là độ dài suffix và
$\mathcal{L}$ là hàm mất mát được thiết kế để hướng mô hình đến một dạng
đầu ra nhất định.

Công trình Universal and Transferable Adversarial Attacks on Aligned
Language Models đề xuất cách tối ưu adversarial suffix bằng kết hợp tìm
kiếm greedy và thông tin gradient, thường được gọi là Greedy Coordinate
Gradient (GCG) \cite{zou2023universal}. Điểm đáng chú ý của hướng này là
suffix được tối ưu trên một tập yêu cầu và mô hình có thể chuyển giao một
phần sang các mô hình khác.

White-box optimization có ý nghĩa quan trọng về mặt nghiên cứu vì nó cho
thấy alignment không chỉ có thể bị phá vỡ bằng prompt thủ công mà còn có
thể bị stress test bằng các phương pháp tối ưu có hệ thống. Tuy nhiên,
thiết lập white-box thường ít phù hợp hơn với các dịch vụ thương mại đóng,
nơi attacker không có quyền truy cập gradient hoặc tham số mô hình.

\subsection{Black-box automated jailbreak}
\label{subsec:black-box-automated-jailbreak}

Trong thiết lập black-box, attacker chỉ có thể gửi prompt và quan sát phản
hồi. Không có gradient trực tiếp để tối ưu đầu vào. Vì vậy, black-box
automated jailbreak thường sử dụng một vòng lặp gồm ba thành phần: một
bộ sinh prompt, một target model và một bộ đánh giá mức độ thành công.

Quy trình tổng quát có thể mô tả như sau:

\begin{enumerate}[leftmargin=1.2cm]
    \item sinh một candidate prompt dựa trên mục tiêu kiểm thử;
    \item gửi candidate prompt đến target model;
    \item đánh giá phản hồi của target model;
    \item cập nhật candidate prompt dựa trên kết quả đánh giá;
    \item lặp lại cho đến khi đạt tiêu chí dừng hoặc hết ngân sách truy vấn.
\end{enumerate}

PAIR là một ví dụ tiêu biểu của black-box automated jailbreak
\cite{chao2023pair}. Phương pháp này sử dụng một attacker LLM để tự động
sinh và điều chỉnh prompt qua nhiều vòng tương tác với target model. Thay
vì dựa hoàn toàn vào prompt thủ công, PAIR xem jailbreak như một quá trình
refinement lặp lại dựa trên phản hồi quan sát được.

TAP mở rộng hướng này bằng cấu trúc tìm kiếm dạng cây và cơ chế pruning
\cite{mehrotra2023tap}. Thay vì chỉ theo một nhánh refinement, TAP duy trì
nhiều candidate prompts, đánh giá sơ bộ và loại bỏ những nhánh ít có khả
năng thành công trước khi gửi đến target model. Cách này giúp giảm số truy
vấn không cần thiết và cải thiện hiệu quả tìm kiếm trong black-box setting.

Các phương pháp black-box automated jailbreak phản ánh bối cảnh thực tế
hơn white-box attack, vì nhiều hệ thống triển khai chỉ cho phép người dùng
gửi truy vấn và nhận phản hồi. Tuy nhiên, chúng cũng đặt ra thách thức lớn
cho evaluation: kết quả phụ thuộc vào ngân sách truy vấn, chất lượng của
attacker model, chất lượng của judge và chính sách dừng.

\subsection{Transferability}
\label{subsec:jailbreak-transferability}

Transferability là khả năng một jailbreak được tạo trên một mô hình hoặc
một tập yêu cầu vẫn có hiệu quả trên mô hình hoặc yêu cầu khác. Đây là một
thuộc tính quan trọng vì attacker không phải lúc nào cũng có quyền truy cập
vào target model thật. Thay vào đó, attacker có thể tối ưu attack trên một
surrogate model rồi thử chuyển sang target model.

Có thể phân biệt ba dạng transferability:

\begin{enumerate}[leftmargin=1.2cm]
    \item \textbf{Cross-prompt transfer:}
    một prompt hoặc suffix có hiệu quả trên nhiều yêu cầu khác nhau;

    \item \textbf{Cross-model transfer:}
    attack được tạo trên một mô hình nhưng vẫn có tác dụng trên mô hình
    khác;

    \item \textbf{Cross-defense transfer:}
    attack vẫn có hiệu quả khi hệ thống bổ sung một số lớp phòng thủ.
\end{enumerate}

Về mặt nghiên cứu, transferability cho thấy các mô hình khác nhau có thể
chia sẻ những điểm yếu tương tự trong cách học hành vi từ chối. Nếu một
attack chỉ hoạt động trên một target duy nhất, nó có thể phản ánh đặc điểm
riêng của mô hình đó. Ngược lại, nếu attack chuyển giao qua nhiều mô hình,
nó gợi ý rằng vấn đề có thể nằm ở cơ chế alignment hoặc dữ liệu căn chỉnh
chung.

Transferability cũng làm tăng mức độ nghiêm trọng của jailbreak. Một attack
có tính chuyển giao cao có thể được phát triển trên mô hình mở hoặc mô hình
nhỏ hơn, sau đó áp dụng vào hệ thống đóng mà attacker không hiểu đầy đủ.
Vì vậy, benchmark jailbreak thường cần đánh giá cả model-specific success
và transfer success.

\subsection{Adaptive attacks}
\label{subsec:adaptive-attacks}

Static attack sử dụng một prompt hoặc chiến lược cố định, không thay đổi
theo defense. Ngược lại, adaptive attack giả định attacker biết hoặc suy
đoán cơ chế phòng thủ và điều chỉnh attack để vượt qua chính cơ chế đó.
Sự khác biệt này rất quan trọng khi đánh giá độ bền vững của defense.

Một defense có thể hoạt động tốt trước static attacks nhưng suy giảm mạnh
khi attacker thích nghi. Ví dụ, nếu hệ thống chỉ chặn một số mẫu bề mặt,
attacker có thể thay đổi cách diễn đạt. Nếu hệ thống sử dụng judge hoặc
guard model, attacker có thể tối ưu prompt để vừa vượt qua guard vừa làm
target model phản hồi không an toàn. Nếu hệ thống dùng nhiều lần kiểm tra
ở thời điểm suy luận, attacker có thể tìm những đầu vào tạo phản hồi ổn
định hơn qua các biến thể.

Có thể biểu diễn adaptive attack như một quá trình tối ưu phụ thuộc vào
defense $D$:

\begin{equation}
    x^{*}
    =
    \arg\max_{x' \in \mathcal{X}}
    A
    \left(
        D(F_{\theta}(x'))
    \right),
    \label{eq:adaptive-attack-objective}
\end{equation}

trong đó $A$ là hàm đánh giá mức độ thành công của attack. Công thức này
chỉ mang tính khái quát, nhưng cho thấy attacker không chỉ tối ưu đối với
mô hình $F_{\theta}$ mà còn tối ưu đối với toàn bộ pipeline có defense.

Adaptive attacks làm rõ một nguyên tắc quan trọng: đánh giá defense chỉ
trên một tập attack cố định thường chưa đủ. Một defense mạnh cần được kiểm
tra trong điều kiện attacker có khả năng quan sát, phản hồi và thay đổi
chiến lược. Đây là lý do các chương sau sẽ phân tích defense không chỉ theo
Attack Success Rate, mà còn theo over-refusal, chi phí suy luận và khả năng
chống lại attacker thích nghi.

\subsection{Tổng kết chương}
\label{subsec:jailbreak-attacks-summary}

Chương này đã hệ thống hóa các nhóm jailbreak attacks theo nhiều chiều:
manual và template-based attacks, obfuscation và encoding attacks,
multi-turn attacks, white-box optimization, black-box automated jailbreak,
transferability và adaptive attacks. Các nhóm này khác nhau về giả định
quyền truy cập, chi phí truy vấn, mức tự động hóa và khả năng chuyển giao.

Một điểm chung giữa các phương pháp là chúng đều khai thác khoảng cách giữa
mục tiêu safety alignment và sự đa dạng gần như không giới hạn của đầu vào
ngôn ngữ tự nhiên. Manual attacks cho thấy mô hình có thể bị ảnh hưởng bởi
cách diễn đạt và bối cảnh. Optimization-based attacks cho thấy có thể tìm
đầu vào bất thường bằng thuật toán. Black-box automated attacks cho thấy
quá trình jailbreak có thể được tự động hóa ngay cả khi không truy cập
tham số mô hình. Multi-turn và adaptive attacks tiếp tục mở rộng bề mặt rủi
ro khi hệ thống có trạng thái, bộ nhớ, công cụ hoặc defense phức tạp.

Các kết quả này dẫn đến nhu cầu nghiên cứu phòng thủ theo nhiều lớp. Chương
tiếp theo phân tích các nhóm jailbreak defenses, bao gồm input filtering,
adversarial training, inference-time defenses, output moderation, guard
models và defense-in-depth.
